\documentclass[11pt]{article}
\usepackage{hyperref}


\setlength{\oddsidemargin}{0in}
\setlength{\evensidemargin}{0in}
\setlength{\topmargin}{-0.5in}
\setlength{\textwidth}{6.5in}
\setlength{\textheight}{9.0in}

\begin{document}
	
	\begin{center}
		{\huge \textbf{Executive Summary: \newline NYC Libraries Service Zones}} \\
		\vspace{2mm}
		\textit{Spatial Analysis and Jurisdictional Mapping}
		\vspace{10mm}
	\end{center}
	
	\noindent \textbf{Overview} \\
	This map provides a spatial analysis of New York City’s three independent library systems: \textbf{The New York Public Library (NYPL)}, the \textbf{Brooklyn Public Library (BPL)}, and the \textbf{Queens Public Library (QPL)}. By transforming point-based branch locations into interactive polygons, the visualization defines the service zones for every public library across the five boroughs.
	
	\vspace{5mm}
	\noindent {\large \textbf{What This Map Tells Us}}
	\begin{itemize}
		\item \textbf{Territorial Dominance \& Jurisdictions:} The map clearly delineates the administrative boundaries of the three systems. It highlights how the NYPL manages a vast, multi-borough territory, while the BPL and QPL maintain high-density, borough-specific networks.
		
		\item \textbf{Service Density Hotspots:} Small polygons in areas like Upper Manhattan indicate high concentration and walkable access. Conversely, larger polygons in eastern Queens or southern Staten Island reveal ``library sprawl'' where residents travel significantly further.
		
		\item \textbf{Strategic Geographic Reach:} By filtering data by system, users can see how each organization positions resources to cover diverse demographics, from high-traffic urban centers to suburban stretches.
	\end{itemize}
	
	\vspace{5mm}
	\noindent {\large \textbf{Why This Information is Compelling}}
	\begin{itemize}
		\item \textbf{Visualizing the ``Invisible'' Border} \\
		While most New Yorkers know their borough, they may not realize where one library system's responsibility ends and another's begins. This map makes those administrative borders visible and interactive.
		
		\item  \textbf{Equity and Access Analysis} \\
		For urban planners, the map serves as a tool for equity. It raises questions about resource allocation---specifically whether large polygons in outlying areas represent a need for more mobile library services or new physical branches.
		
	 	\item  \textbf{Visualizing Intellectual Infrastructure} \\
		This high-contrast interface serves as a functional blueprint of the city's knowledge network. By isolating spatial data against a dark background, the map emphasizes the rigid ``grid'' of urban service delivery, transforming standard municipal records into a clear narrative about how New York City distributes its intellectual resources.
	\end{itemize}
	
	\vspace{5mm}
	\noindent  {\large \textbf{Data Source}} \\
	The underlying geographic and demographic data used in this analysis was sourced from \newline 
	\textbf{NYC Open Data}: 
	\href{https://data.cityofnewyork.us/Business/Library/p4pf-fyc4}{NYC Library Map}.
	
\end{document}